\chapter{Background and Problem Formulation}
\label{cap:background}

This chapter establishes the agricultural navigation context and formulates the vision-based local guidance problem addressed in this thesis.

\section{Agricultural autonomous navigation}
\label{sec:agricultural_navigation}

\subsection{Navigation tasks and requirements}
\label{subsec:navigation_tasks}

Agricultural robots perform tasks such as monitoring, spraying, harvesting, and weeding in structured crop environments~\cite{bai2023vision}. Navigation between crop rows requires the robot to follow a path aligned with the row structure while avoiding plants and obstacles. The guidance reference must be computed from local sensing, as global GNSS positioning may be obstructed by crop canopies or greenhouse structures~\cite{bai2023vision,navone2025gpsfree}.

Vision-based methods identify crop rows, inter-row paths, or crop boundaries from camera images and convert these observations into steering commands~\cite{bai2023vision,silva2023rowdetection}. The task is local: the robot maintains its position relative to the visible row structure, not a global map.

\subsection{Environmental challenges}
\label{subsec:environmental_challenges}

Field conditions vary with illumination, shadows, weather, growth stage, and weed presence~\cite{kneip2020cropedge,silva2023rowdetection}. Plants may occlude parts of the path, and gaps or irregular spacing disrupt row structure. These changes affect both the appearance of image features and the reliability of geometric references extracted from them.

Segmentation models trained on one set of conditions must generalize to new lighting, crops, or row configurations. Geometric extractors must handle incomplete or ambiguous corridor observations without assuming perfect segmentation~\cite{li2022multiroi}.

\section{Vision-based local guidance}
\label{sec:vision_guidance}

\subsection{Sensing modalities}
\label{subsec:sensing_modalities}

Camera-based perception provides appearance information for distinguishing plants, soil, and paths. Monocular RGB cameras are simpler and less expensive than LiDAR or stereo systems, but do not directly measure depth~\cite{bai2023vision}. While RGB-D cameras provide both appearance and depth measurements, their performance depends on field conditions and hardware cost constraints.

This thesis uses monocular RGB sensing for local guidance. It does not reconstruct three-dimensional field structure or maintain a global map.

\subsection{Image-space reference extraction}
\label{subsec:image_space_reference}

The guidance pipeline converts image observations into a reference suitable for control. Depending on the method, this may be a horizontal position, a fitted line, or a curve~\cite{chen2021navigation,li2022multiroi,mendez2023vineyards}. Image-space references describe the path geometry in pixel coordinates; they must be interpreted by the controller to generate velocity commands.

Straight-line models assume locally linear boundaries and produce simple references for proportional control~\cite{mendez2023vineyards}. They do not capture distant curves or complex row patterns, but are computationally efficient and suitable for short prediction horizons.

\placeholderfigure{fig:field_corridor_geometry}%
{Agricultural inter-row environment, vehicle frames, and geometric error definition}%
{Schema geometrico dei sistemi di coordinate: terna fissa $\{W\}$, terna robot $\{R\}$, terna ottica della telecamera $\{C\}$ e piano immagine $\{I\}$. Rappresentazione grafica dell'offset laterale $e_y$ rispetto alla linea centrale del corridoio e dell'errore di orientamento $e_\theta$.}%
{5.5cm}

\todobox{Modello cinematico 4WS e formulazione analitica degli errori}{%
Integrare la formalizzazione matematica del veicolo e del controllo:
\begin{itemize}
  \item Modello cinematico per piattaforma 4WS a sterzatura simmetrica in controfase ($\delta_f = -\delta_r$) e relazione cinematica: $\dot{x} = v \cos\theta$, $\dot{y} = v \sin\theta$, $\dot{\theta} = \omega = \frac{2v\tan\delta}{L}$.
  \item Trasformazione rigida e proiettiva tra coordinate pixel $(u, v)$ nel piano immagine e coordinate metriche nel frame robot $\{R\}$.
  \item Definizione rigorosa dell'offset laterale $e_y$ e dell'heading error $e_\theta$ calcolati rispetto alla retta di mezzeria stimata.
  \item Formulazione della legge di guida proporzionale / Stanley per la determinazione della velocit\`a angolare $\omega = K_p e_y + K_\theta e_\theta$.
\end{itemize}
}

\section{The segmentation-to-control pipeline}
\label{sec:pipeline_structure}

\placeholderfigure{fig:overall_system_pipeline}%
{Complete vision-based local guidance pipeline from camera stream to actuator commands}%
{Schema a blocchi dell'architettura di navigazione: flusso continuo dall'acquisizione delle immagini monoculari, attraverso la rete di segmentazione semantica, l'estrazione geometrica dei bordi del corridoio, il calcolo degli errori di guida e la generazione dei comandi di velocit\`a $(v, \omega)$.}%
{4.5cm}

\subsection{Perception stage}
\label{subsec:perception_stage}

The perception stage assigns semantic labels to image pixels. For inter-row navigation, target classes typically comprise crop plants, the traversable path, covering material such as plastic mulch, and background~\cite{li2022multiroi}. The definition of output classes directly determines what the geometric stage can extract: a binary vegetation mask identifies crop rows, whereas a path mask isolates the drivable space between them~\cite{silva2023rowdetection,navone2025gpsfree}.

Segmentation models are trained on annotated field images. While transfer learning from pretrained backbones reduces training duration and improves feature generalization, the classification head must be tailored to the specific target crop classes and deployed at low latency on embedded hardware~\cite{navone2025gpsfree}.

\subsection{Geometric extraction stage}
\label{subsec:geometric_stage}

The geometric stage reduces the dense segmentation mask to a low-dimensional guidance reference for robot steering. Existing methods differ in their spatial search strategy and fitted geometric model: some aggregate mask pixels across image columns~\cite{mendez2023vineyards,navone2025gpsfree}, whereas others extract local intervals from horizontal strips~\cite{li2022multiroi} or search for line candidates within bounded angular regions~\cite{chen2021navigation}.

The extraction algorithm must tolerate segmentation noise, missing plants, and cast shadows without producing erratic steering targets. Validation checks can reject estimates that fail geometric consistency thresholds, requiring the controller to handle transient frames where no valid reference is recovered~\cite{kneip2020cropedge}.

\subsection{Control integration}
\label{subsec:control_integration}

The controller converts the extracted geometric reference into forward velocity and steering commands. Proportional and look-ahead laws regulate steering using the lateral displacement and heading orientation errors relative to the image center~\cite{mendez2023vineyards}.

In field operations, raw error signals extracted from semantic segmentation exhibit high-frequency jitter and occasional impulse disturbances caused by irregular foliage contours or transient detection dropouts. Directly feeding unfiltered errors into the steering loop induces actuator chattering and abrupt steering adjustments that compromise vehicle stability. Consequently, the control interface incorporates error conditioning structured across three complementary mechanisms:
\begin{enumerate}
    \item \textbf{High-frequency smoothing:} a recursive low-pass filter attenuates persistent image-space tracking jitter without introducing excessive phase delay during row transitions;
    \item \textbf{Impulse outlier rejection:} large step variations between consecutive frames are evaluated over temporal windows to isolate single-frame false detections while preserving genuine trajectory adjustments;
    \item \textbf{Temporal hold and deadband:} a zero-order hold maintains command continuity across short-term perception dropouts, while a symmetric deadband around the corridor center prevents steering hunting.
\end{enumerate}
The mathematical formulation and parameter tuning of this conditioning pipeline on the Demetra robotic platform are detailed in Section~\ref{subsubsec:control_node} of Chapter~\ref{cap:platform}.

Closed-loop field evaluation measures the robot's ability to maintain row alignment under operational conditions. High segmentation accuracy and accurate geometric fitting are necessary prerequisites, but stable inter-row tracking ultimately depends on real-time execution and disturbance rejection over continuous runs~\cite{mendez2023vineyards,silva2023rowdetection,navone2025gpsfree}.
